\section{Request and receipt semantics}

\subsection{Canonical request}

\begin{definition}[Execution request]
An execution request is the canonical tuple
\[
q=(p,c,i,e,v,t,m,s,b),
\]
where $p$ is the protocol digest, $c$ the code digest, $i$ the ordinary input
manifest, $e$ the scientific-evidence manifest, $v$ the environment and tool
manifest, $t$ the requested terminal schema, $m$ the required executor/model
identity policy, $s$ the seed and session policy, and $b$ the resource and access
budget. Its request identifier is $H(\operatorname{canon}(q))$.
\end{definition}

Canonicalization is part of the protocol, not an implementation detail. Every
field that changes observable execution or admissibility must either occur in $q$
or be prohibited. The request identifier therefore binds declared semantics, not
only a prompt string. The scientific-evidence manifest is distinct from ordinary
inputs because source bytes, credentials, caches and auxiliary configuration have
different authority.

\subsection{Result receipt}

\begin{definition}[Result receipt]
A result receipt is
\[
r=(h_q,h_o,z,u,k,g,w,\tau,\sigma),
\]
where $h_q$ is the request identifier, $h_o$ the output-manifest digest, $z$ a
typed terminal, $u$ an authenticated executor subject, $k$ provider/model/checkpoint
identity when applicable, $g$ the realized code/tool/environment configuration,
$w$ the evidence-lineage graph, $\tau$ timing and epoch data, and $\sigma$ an
attestation over the canonical receipt payload.
\end{definition}

The terminal $z$ belongs to a closed sum type:
\[
z\in\{\textsc{completed},\textsc{valid-negative},\cannotcheck,
       \failure,\textsc{invalid-result}\}.
\]
A host or capability failure is therefore not encoded as a numerical outcome or a
successful empty result. A valid negative is a result of the scientific protocol;
an execution failure is a result of the computing environment. Conflating them is
a type error.

\subsection{Verification predicates}

For a request $q$, receipt $r$ and trust policy $T$, define:

\begin{align*}
S(q,r) &:\text{ schema and canonicalization are valid},\\
D(q,r) &:\text{ request, output and manifest digests match},\\
I(q,r,T) &:\text{ executor/model identities satisfy }T,\\
C(q,r,T) &:\text{ attestation chain and issuer scope satisfy }T,\\
E(q,r) &:\text{ evidence lineage is complete for the claimed level},\\
Z(r) &:\text{ terminal is admissible for the requested scientific endpoint}.
\end{align*}

The mechanical admission predicate is
\[
M(q,r,T)=S\land D\land I\land C\land E.
\]
Scientific admission additionally requires a frozen adjudication rule $F$:
\[
\operatorname{Admit}(q,r,T,F)=M(q,r,T)\land Z(r)\land F(q,r).
\]
No missing predicate is coerced to true. Missing or unverifiable identity,
lineage, trust or protocol data produce \cannotcheck{} for the associated claim.

